
\documentclass[11pt,spanish]{article}

% =========================================================
% PLANTILLA BASE PARA INFORMES ACADÉMICOS / TÉCNICOS
% =========================================================

% --- Idioma y codificación ---
\usepackage[utf8]{inputenc}
\usepackage[spanish,es-tabla]{babel}
\usepackage[T1]{fontenc}
\usepackage{textcomp}
\usepackage{newunicodechar}
\newunicodechar{₡}{\textcolonmonetary}

% --- Márgenes / papel ---
\usepackage[
    left=2.5cm,
    right=2.5cm,
    top=2.5cm,
    bottom=2.5cm,
    headsep=0.5cm,
    footskip=0.8cm
]{geometry}

% --- Matemáticas ---
\usepackage{amsmath}

% --- Tablas, enlaces y elementos visuales ---
\usepackage{array}
\usepackage{tabularx}
\usepackage{booktabs}
\usepackage{longtable}
\usepackage{graphicx}
\usepackage{float}
\usepackage{xcolor}
\usepackage{tikz}
\usetikzlibrary{positioning,shapes.geometric,arrows.meta,calc}
\usepackage{enumitem}
\usepackage{ragged2e}
\usepackage[hidelinks]{hyperref}

% --- Encabezado y pie de página ---
\usepackage{fancyhdr}
\setlength{\headheight}{14pt}
\usepackage{lastpage}

% --- Interlineado ---
\linespread{1.2}
\sloppy

% =========================================================
% DATOS DEL DOCUMENTO
% =========================================================
\newcommand{\Institucion}{Instituto Tecnológico de Costa Rica}
\newcommand{\Campus}{Campus Tecnológico Central Cartago}
\newcommand{\DocumentoTitulo}{Plan de Adquisiciones}
\newcommand{\Curso}{Administración de proyectos}
\newcommand{\Profesor}{Alicia Marcela Salazar Hernández}
\newcommand{\FechaDocumento}{17 de junio de 2026}
\newcommand{\PeriodoAcademico}{I Semestre}
\newcommand{\AnioDocumento}{2026}
\newcommand{\AutoresPortada}{%
    \begin{tabular}{l@{\hspace{1.5cm}}r}
        Mauricio González Prendas & 2024143009 \\
        Isaac Jiménez Blanco & 2024202804 \\
        Tamara Robles Camacho & 2024099342 \\
    \end{tabular}%
}
\newcommand{\DocHeaderLeft}{}
\newcommand{\DocHeaderRight}{Plan de Adquisiciones}
\newcommand{\DocFooterRight}{Página~\thepage\ de~\pageref{LastPage}}

% Metadatos PDF.
\title{\DocumentoTitulo}
\author{\AutoresPortada}
\date{\FechaDocumento}

% =========================================================
% CONFIGURACIÓN DE TABLAS Y UTILIDADES
% =========================================================
\newcolumntype{Y}{>{\RaggedRight\arraybackslash}X}
\newcolumntype{L}[1]{>{\RaggedRight\arraybackslash}p{#1}}
\renewcommand{\arraystretch}{1.22}
\setlength{\LTpre}{0.5em}
\setlength{\LTpost}{0.5em}

% Imagen segura: si el archivo no existe, muestra un recuadro de marcador.
\newcommand{\SafeImage}[2][0.92\textwidth]{%
    \IfFileExists{#2}{%
        \includegraphics[width=#1]{#2}%
    }{%
        \fbox{%
            \begin{minipage}[c][4cm][c]{#1}
            \centering
            Imagen pendiente:\\
            \texttt{#2}
            \end{minipage}%
        }%
    }%
}

% Figura reutilizable.
\newcommand{\EvidenceFigure}[3]{%
    \begin{figure}[H]
    \centering
    \SafeImage{#1}
    \caption{#2}
    \label{#3}
    \end{figure}%
}

% =========================================================
% ENCABEZADO Y PIE
% =========================================================
\pagestyle{fancy}
\fancyhf{}
\fancyhead[L]{\DocHeaderLeft}
\fancyhead[R]{\DocHeaderRight}
\renewcommand{\headrulewidth}{0.4pt}
\renewcommand{\footrulewidth}{0.4pt}
\fancyfoot[R]{\DocFooterRight}

% Estilo equivalente para páginas horizontales.
\fancypagestyle{widefancy}{%
    \fancyhf{}%
    \fancyhead[L]{\DocHeaderLeft}%
    \fancyhead[R]{\DocHeaderRight}%
    \renewcommand{\headrulewidth}{0.4pt}%
    \renewcommand{\footrulewidth}{0.4pt}%
    \fancyfoot[R]{\DocFooterRight}%
}

% =========================================================
% PÁGINAS HORIZONTALES PARA TABLAS ANCHAS
% =========================================================
\makeatletter
\newcommand{\SyncPageBuilderDimensions}{%
    \global\vsize=\textheight
    \global\@colroom=\textheight
    \global\@colht=\textheight
}
\makeatother

\newenvironment{widepage}{%
    \clearpage
    \hypersetup{pageanchor=false}%
    \paperwidth=297mm
    \paperheight=210mm
    \pdfpagewidth=297mm
    \pdfpageheight=210mm
    \newgeometry{left=2.5cm,right=2.5cm,top=2.5cm,bottom=2.5cm,headsep=0.5cm,footskip=0.8cm}%
    \setlength{\textwidth}{243mm}%
    \setlength{\textheight}{156mm}%
    \setlength{\linewidth}{\textwidth}%
    \setlength{\columnwidth}{\textwidth}%
    \hsize=\textwidth
    \setlength{\headwidth}{\textwidth}%
    \SyncPageBuilderDimensions
    \pagestyle{widefancy}%
}{%
    \clearpage
    \paperwidth=210mm
    \paperheight=297mm
    \pdfpagewidth=210mm
    \pdfpageheight=297mm
    \restoregeometry
    \SyncPageBuilderDimensions
    \hypersetup{pageanchor=true}%
    \pagestyle{fancy}%
}

% =========================================================
% PORTADA
% =========================================================
\newcommand{\MakeCover}{%
\begin{titlepage}
\thispagestyle{empty}
\centering

{\bfseries \huge \Institucion \par}
\vspace{0.25cm}
{\LARGE \Campus \par}

\vfill

{\huge \DocumentoTitulo \par}

\vfill

{\bfseries \LARGE Profesora \par}
{\Large \Profesor \par}

\vfill

{\bfseries \LARGE Estudiantes \par}
{\Large \AutoresPortada \par}

\vfill

{\bfseries\LARGE Fecha \par}
{\Large \FechaDocumento \par}

\vfill

{\LARGE \PeriodoAcademico \par}
{\LARGE \AnioDocumento \par}

\end{titlepage}%
}

% =========================================================
% DOCUMENTO
% =========================================================
\begin{document}
\hypersetup{pageanchor=false}

% =========================
% Portada
% =========================
\MakeCover

% =========================
% Índice
% =========================
\pagenumbering{roman}
\pagestyle{empty}
\addtocontents{toc}{\protect\thispagestyle{empty}}
\tableofcontents
\clearpage
\pagenumbering{arabic}
\hypersetup{pageanchor=true}
\pagestyle{fancy}

% =========================
% Contenido
% =========================

\section{Información general del documento}

\begin{table}[H]
\centering
\begin{tabularx}{\textwidth}{p{5cm} Y}
\toprule
\textbf{Nombre del proyecto} & Plataforma Universitaria Integrada \\
\textbf{Organización} & Universidad Tecnológica La Mejor \\
\textbf{Documento} & \DocumentoTitulo \\
\textbf{Directora del proyecto} & Tamara Robles Camacho \\
\textbf{Fecha} & \FechaDocumento \\
\bottomrule
\end{tabularx}
\end{table}

\section{Introducción}

El presente documento corresponde al Plan de Adquisiciones del proyecto Plataforma Universitaria Integrada, desarrollado para la Universidad Tecnológica La Mejor. Este plan forma parte de la planificación del proyecto y tiene como finalidad definir los recursos externos, herramientas, servicios, materiales y apoyos operativos que deben considerarse para ejecutar el trabajo de manera ordenada durante el periodo planificado de tres meses.

El proyecto consiste en la elaboración de un prototipo web navegable y de un conjunto de documentos de administración de proyectos. Debido a que el alcance aprobado excluye el desarrollo de backend funcional, base de datos real, autenticación institucional, pasarela de pagos real y despliegue productivo obligatorio, las adquisiciones no se orientan a comprar una solución empresarial completa, sino a sostener la elaboración del prototipo, la documentación, la comunicación del equipo, el cronograma, las pruebas y las evidencias finales.

El plan se relaciona directamente con el cronograma elaborado en Microsoft Project, el plan de recursos, la estimación de costos y el plan de calidad. En consecuencia, las adquisiciones se manejan de manera controlada para evitar gastos innecesarios y asegurar que el presupuesto total del proyecto se mantenga en ₡17,625,200. El enfoque utilizado prioriza herramientas disponibles, licencias temporales, servicios de bajo costo y recursos ya existentes, ya que el proyecto tiene naturaleza académica y no requiere una contratación institucional completa.

\section{Propósito del plan de adquisiciones}

El propósito de este documento es establecer qué elementos deben adquirirse, contratarse, utilizarse o reservarse para apoyar el desarrollo del proyecto. Asimismo, permite definir criterios básicos para seleccionar recursos, asignar responsables, controlar costos y mantener trazabilidad entre las necesidades del proyecto y el presupuesto aprobado.

Este plan busca cumplir los siguientes objetivos:

\begin{itemize}[leftmargin=*, label={\textbullet}]
    \item Identificar los recursos externos o complementarios requeridos por el proyecto.
    \item Diferenciar entre recursos ya disponibles y recursos que podrían representar un costo.
    \item Establecer criterios simples para seleccionar herramientas, servicios y materiales.
    \item Relacionar las adquisiciones con el cronograma, el presupuesto y el alcance aprobado.
    \item Evitar compras innecesarias o elementos fuera del alcance del prototipo académico.
    \item Asignar responsables para la revisión, aprobación y seguimiento de cada adquisición.
\end{itemize}

\section{Alcance de las adquisiciones}

Las adquisiciones consideradas en este proyecto se limitan a recursos de apoyo para la gestión, documentación, desarrollo y presentación final. No se contempla la compra de un ERP universitario, contratación de servidores productivos definitivos, licencias empresariales de gran escala ni servicios profesionales externos para desarrollar el prototipo.

\begin{table}[H]
\centering
\caption{Alcance de las adquisiciones}
\label{tab:alcance-adquisiciones}
{\small
\begin{tabularx}{\textwidth}{p{4cm} Y}
\toprule
\textbf{Tipo de recurso} & \textbf{Alcance definido} \\
\midrule
Software y herramientas & Licencias temporales o herramientas de apoyo para cronograma, documentación, diseño, desarrollo y control de versiones. \\
Hardware y materiales & Recursos físicos mínimos, almacenamiento externo y periféricos complementarios si fueran requeridos. \\
Servicios operativos & Conexión a internet, impresiones, transporte a reuniones o servicios básicos de apoyo al trabajo del equipo. \\
Hosting o infraestructura & Uso de alternativas gratuitas o de bajo costo para respaldo, repositorio y posible demostración del prototipo. \\
Contingencia & Reserva para atender imprevistos relacionados con herramientas, materiales, conectividad o necesidades menores de ejecución. \\
\bottomrule
\end{tabularx}
}
\end{table}


\section{Análisis Make vs Buy}

Para este proyecto se aplica un análisis Make vs Buy con el fin de definir qué elementos serán desarrollados por el equipo interno y cuáles recursos deberán adquirirse o utilizarse como apoyo externo. Debido a que el alcance corresponde a un prototipo académico Front-End navegable y no a un sistema productivo completo, la estrategia principal será desarrollar internamente los entregables técnicos y documentales, mientras que las adquisiciones se limitarán a herramientas, servicios y materiales de apoyo.

\begin{table}[H]
\centering
\caption{Análisis Make vs Buy del proyecto}
\label{tab:make-buy}
{\small
\begin{tabularx}{\textwidth}{p{4cm} p{3cm} Y}
\toprule
\textbf{Elemento} & \textbf{Decisión} & \textbf{Justificación} \\
\midrule
Prototipo Front-End navegable & Desarrollar internamente & El equipo cuenta con desarrolladores asignados y el alcance exige construir una interfaz web navegable, no comprar un sistema comercial. \\
Documentación del proyecto & Desarrollar internamente & Los entregables forman parte de la evaluación académica y deben ser elaborados por el equipo según los contenidos vistos en clase. \\
Cronograma, WBS, RACI, riesgos, calidad y cambios & Desarrollar internamente & Estos documentos representan la aplicación práctica de la administración de proyectos y deben mantener coherencia con el alcance aprobado. \\
Microsoft Project y herramientas de documentación & Adquirir o utilizar como apoyo & Se requieren para elaborar cronograma, evidencias, documentos y archivos finales, pero no sustituyen el trabajo del equipo. \\
Hosting o infraestructura de demostración & Utilizar apoyo externo básico & Puede usarse una alternativa gratuita o de bajo costo únicamente para mostrar evidencia del prototipo, sin implicar despliegue productivo. \\
ERP universitario, backend, base de datos real o pasarela de pagos & No adquirir & Estos elementos están fuera del alcance aprobado y aumentarían innecesariamente el costo, el tiempo y la complejidad del proyecto. \\
\bottomrule
\end{tabularx}
}
\end{table}

Con base en este análisis, la decisión general del proyecto es desarrollar internamente el prototipo y la documentación, y adquirir únicamente recursos complementarios de bajo costo. Esta estrategia permite mantener el presupuesto total de ₡17,625,200 y evita compras innecesarias que podrían generar alcance creciente o contradicciones con la definición del alcance.


\section{Criterios de selección}

Para seleccionar cualquier recurso o servicio se utilizarán criterios simples y coherentes con la naturaleza del proyecto. Estos criterios permiten mantener el control financiero y evitar que una adquisición altere el alcance o genere dependencia innecesaria.

\begin{table}[H]
\centering
\caption{Criterios de selección de adquisiciones}
\label{tab:criterios-adquisiciones}
{\small
\begin{tabularx}{\textwidth}{p{4cm} Y}
\toprule
\textbf{Criterio} & \textbf{Aplicación en el proyecto} \\
\midrule
Costo razonable & La adquisición debe mantenerse dentro del presupuesto aprobado y preferiblemente utilizar opciones gratuitas o temporales. \\
Necesidad para el alcance & Solo se aceptan recursos relacionados con documentación, cronograma, prototipo Front-End, pruebas o evidencias finales. \\
Facilidad de uso & La herramienta o servicio debe poder ser utilizado por el equipo sin capacitación extensa. \\
Disponibilidad inmediata & Se priorizan recursos que puedan ser utilizados durante el ciclo del proyecto sin procesos largos de contratación. \\
Compatibilidad técnica & Las herramientas deben ser compatibles con tecnologías web, repositorio de versiones, PDF, Microsoft Project o documentación académica. \\
Trazabilidad & Las decisiones relevantes deben quedar documentadas en el repositorio, minuta, correo o archivo de seguimiento cuando corresponda. \\
\bottomrule
\end{tabularx}
}
\end{table}

\section{Plan de adquisiciones del proyecto}

La siguiente tabla presenta el plan principal de adquisiciones. Los montos se mantienen alineados con la estimación PERT y con el presupuesto actualizado del proyecto.

\begin{table}[H]
\centering
\caption{Plan de adquisiciones}
\label{tab:plan-adquisiciones}
{\scriptsize
\begin{tabularx}{\textwidth}{p{2.6cm} p{3.3cm} Y p{2.4cm} p{2.1cm}}
\toprule
\textbf{Categoría} & \textbf{Recurso} & \textbf{Uso previsto} & \textbf{Responsable} & \textbf{Costo estimado} \\
\midrule
Software & Microsoft Project & Elaborar cronograma, tareas, dependencias, hitos, recursos, costos y ruta crítica. & Isaac Jiménez Blanco & ₡66,667 \\
Software & Office 365 o herramientas equivalentes & Elaboración y revisión de documentos, tablas, evidencias y archivos finales en PDF. & Equipo interno & ₡131,667 \\
Software & Herramientas premium opcionales & Apoyo visual o utilidades menores en caso de requerirse durante la documentación o presentación. & Tamara Robles Camacho & ₡23,333 \\
Hardware & Almacenamiento externo & Respaldo local de archivos finales, cronograma, evidencias y documentos importantes. & Isaac Jiménez Blanco & ₡19,167 \\
Hardware & Periféricos complementarios & Apoyo operativo mínimo para trabajo local, revisiones o entrega de evidencias. & Equipo interno & ₡14,167 \\
Servicios & Conexión a internet & Comunicación, repositorio, reuniones virtuales, investigación, desarrollo y carga de evidencias. & Cada integrante & ₡359,167 \\
Servicios & Impresiones y material físico & Material de apoyo o evidencias impresas si la entrega académica lo solicita. & Tamara Robles Camacho & ₡10,500 \\
Servicios & Transporte a reuniones presenciales & Traslados mínimos en caso de reuniones de coordinación o revisión presencial. & Equipo interno & ₡41,333 \\
\bottomrule
\end{tabularx}
}
\end{table}

El subtotal de rubros adicionales asociados a software, hardware y servicios corresponde a ₡666,000. Además, el proyecto considera una reserva de contingencia de ₡1,515,867 para atender imprevistos sin modificar el presupuesto total autorizado.

\section{Responsables y aprobación}

Las adquisiciones menores podrán ser revisadas por el equipo interno, pero cualquier gasto que afecte el presupuesto total, el alcance, el cronograma o la calidad deberá ser validado por la directora del proyecto. Si una adquisición representa un cambio relevante o genera una variación presupuestaria, deberá tratarse mediante el proceso formal de control de cambios.

\begin{table}[H]
\centering
\caption{Responsables del plan de adquisiciones}
\label{tab:responsables-adquisiciones}
{\small
\begin{tabularx}{\textwidth}{p{4cm} Y}
\toprule
\textbf{Rol} & \textbf{Responsabilidad} \\
\midrule
Tamara Robles Camacho & Validar necesidades de adquisición, controlar que los gastos se mantengan dentro del presupuesto y aprobar ajustes relevantes. \\
Isaac Jiménez Blanco & Mantener trazabilidad documental de herramientas, costos, evidencias y relación con el cronograma del proyecto. \\
Jorge Abarca Quiros & Apoyar la revisión de costos, adquisiciones menores y coherencia con la estimación PERT. \\
Mauricio González Prendas & Informar necesidades técnicas del desarrollo Front-End, repositorio, pruebas y herramientas relacionadas con los portales asignados. \\
Carlos Sainz Arce & Informar necesidades técnicas del portal financiero, dashboard ejecutivo y pruebas de calidad del prototipo. \\
\bottomrule
\end{tabularx}
}
\end{table}

\section{Gestión y seguimiento de adquisiciones}

El seguimiento de las adquisiciones se realizará durante la planificación, ejecución y cierre del proyecto. Para ello se verificará que los recursos se utilicen únicamente para apoyar las tareas definidas en el cronograma y que no generen nuevas obligaciones fuera del alcance aprobado.

\begin{itemize}[leftmargin=*, label={\textbullet}]
    \item Revisar que las herramientas utilizadas estén disponibles antes de iniciar las tareas críticas.
    \item Confirmar que Microsoft Project permita generar evidencias del cronograma y la ruta crítica.
    \item Mantener respaldos de documentos, código y archivos finales en el repositorio del proyecto.
    \item Registrar cualquier gasto relevante para mantener coherencia con la estimación de costos.
    \item Usar alternativas gratuitas siempre que cumplan con los requisitos del proyecto.
    \item Evitar contrataciones o compras que impliquen backend, base de datos real o operación productiva del sistema.
\end{itemize}

\section{Riesgos asociados a las adquisiciones}

Las adquisiciones del proyecto se relacionan con riesgos moderados debido a que dependen de herramientas tecnológicas y servicios disponibles durante el periodo de trabajo. Los principales riesgos son pérdida de acceso a herramientas, problemas con Microsoft Project, dificultad para exportar evidencias, fallos en el repositorio o inconsistencias entre costos documentados.

Para controlar estos riesgos, el equipo mantendrá copias locales de los archivos importantes, respaldos en GitHub, exportaciones en PDF y revisión cruzada de los montos reportados en documentos como cronograma, plan de calidad, registro de riesgos y estimación de costos.

\section{Conclusión}

El Plan de Adquisiciones de la Plataforma Universitaria Integrada establece una guía simple y suficiente para identificar, controlar y justificar los recursos requeridos durante el proyecto. Debido a que el trabajo se centra en un prototipo académico Front-End y en documentos de gestión, las adquisiciones se enfocan en herramientas de planificación, documentación, desarrollo, conectividad, respaldo y evidencias finales.

La estrategia adoptada permite mantener el presupuesto total de ₡17,625,200, aprovechar recursos ya disponibles y evitar compras innecesarias que puedan alterar el alcance aprobado. Asimismo, el plan facilita la coordinación entre la directora del proyecto, los analistas documentadores y los desarrolladores, de manera que cualquier necesidad adicional pueda revisarse de forma ordenada y, si corresponde, gestionarse mediante control formal de cambios.

\end{document}
